\documentclass[11pt]{article}

% ---- Packages ----
\usepackage[margin=1in]{geometry}
\usepackage{amsmath}
\usepackage{amssymb}
\usepackage{booktabs}
\usepackage{enumitem}
\usepackage{titlesec}
\usepackage[hidelinks]{hyperref}
\usepackage{xcolor}

% ---- Light formatting ----
\definecolor{moss}{HTML}{758C58}
\titleformat{\section}{\large\bfseries\color{moss}}{\thesection}{0.6em}{}
\titleformat{\subsection}{\normalsize\bfseries}{\thesubsection}{0.6em}{}
\setlist{itemsep=2pt, topsep=4pt}
\setlength{\parskip}{0.4em}
\setlength{\parindent}{0pt}

% ---- Title block ----
\title{\vspace{-1.5em}\textbf{Claude Fears the Word ``Loss''}\\[0.3em]
\large Phase 2 Results: Framing, Loss Aversion, and a Violation of Frame Invariance}
\author{J. Blakeslee \\ \small RAND School of Public Policy}
\date{\small AI Sprint Practicum (Course 1012) \\ \small July 2026}

\begin{document}
\maketitle

\hrule
\vspace{0.5em}

\section{Introduction}

This project treats large language models as economic agents and recovers their preference primitives
from their choices. Phase 1 fit a structural discrete-choice model to 1{,}200 controlled lottery
choices from a single frontier model and found something sharp: \texttt{claude-haiku-4-5} prices the
probability of winning and nearly ignores the size of the prize, with an extreme S-shaped probability
weighting function (Prelec $\gamma = 2.69$) and human-plausible curvature once that distortion is
controlled. Phase 1 measured how the model chooses among risky prospects. It left open a different
question: does the model care not just about outcomes, but about how those outcomes are described?

Phase 2 answers that. Expected-utility theory rests on frame invariance: a rational agent evaluates
final states of wealth, so two descriptions of the identical lottery must yield the identical choice.
The classic human violation is Tversky and Kahneman's (1981) Asian disease problem, where reversing
``lives saved'' and ``lives lost'' reverses the majority preference over outcomes that are arithmetically
the same. We port that manipulation to money. The same terminal-wealth lottery is presented under three
frames that a frame-invariant agent must treat identically, and we structurally estimate reference
dependence and loss aversion from the pattern of choices across frames. The model under study is again
\texttt{claude-haiku-4-5-20251001}, a pinned snapshot, queried at temperature 1.0 through forced
tool-choice with an 8\% plain-text validation arm and fresh context per trial.

The headline is that the model is not frame-invariant. Two lotteries with byte-for-byte identical
terminal wealth draw 40\% versus 10\% gamble-taking based only on whether the downside is worded as a
loss. An expected-utility maximizer cannot produce this. The word ``loss'' is doing the work.

\section{The Result in Plain Language}

We asked the same AI model the same kind of money question three different ways, where the three ways
lead to exactly the same outcomes. Imagine you are offered a sure \$50, or a gamble that pays \$120 if
you win and nothing if you lose. That is one version. Now imagine we first hand you \$50 and say it is
yours, then offer you the choice to keep it or gamble: the gamble either raises your \$50 to \$120, or
drops it back to zero. If you do the arithmetic, these are the same decision. Keep the \$50, or take the
gamble that ends with \$120 or \$0. A careful thinker should not care which way we phrased it.

The model cared, and cared a lot. When the downside was described plainly, it took the gamble about
40\% of the time. When the identical downside was described as ``dropping'' from \$50 back to zero, a
loss of the money it had just been handed, it took the gamble only about 10\% of the time. Same money,
same odds, same final outcomes. The only thing that changed was the word ``loss.'' That single change in
wording cut the model's willingness to gamble by three quarters.

This is a well-documented human quirk. People hate losses more than they enjoy equal-sized gains, and
they respond to whether a choice is framed as a gain or a loss even when the outcomes are the same. The
striking thing is that the AI shows the same quirk, and by our best estimate shows it more strongly than
people typically do. When we put a number on how much more the model dislikes a loss than it likes an
equal gain, we get a factor above three. The classic human number is around 2.25. So the model is loss-
averse in the same direction as humans, and by this measure harder.

The behavior was not total. Even in the loss-worded version, the model still gambled a bit more when the
gamble was a better deal, so it was not blindly refusing. It was heavily suppressed, not switched off.
We also ran a simple honesty check: offer the model a sure \$50 against a sure \$70, no gambling
involved. It picked the larger amount every single time, all eighty times. So the model reads how much
money is on the table perfectly well when there is no gamble in the picture. Its trouble is specific to
choices under risk, and here, specific to the word ``loss.''

Put the two phases together and a picture starts to form, carefully. Phase 1 found that when the model
faces a gamble, it fixates on the odds and barely notices the prize size. Phase 2 finds that it also
fixates on loss wording. Yet the honesty check shows it reads dollar amounts fine when no gamble is
present. The tentative reading is that the model's risky choices are steered by a few loud surface
features of the prompt, the probability figure and the loss language, more than by a full calculation of
what it ends up with. That is a hypothesis the two phases jointly support, not a proven fact.

Why does this matter? An agent whose willingness to take a risk can be swung thirty percentage points by
a choice of words is steerable by whoever writes the words. In procurement, negotiation, insurance, or
triage tools, describing an option's downside as a loss becomes a lever over what the agent decides, and
a way to manipulate it in an adversarial setting.

\section{Experimental Design}

Every choice is an independent API call with no memory of prior trials. The instrument is a forced
binary choice between a sure amount \$$s$ and a two-outcome gamble that pays a high amount \$$h$ with
probability $p$ and \$0 otherwise. The manipulation is the frame: the same terminal-wealth lottery is
worded three ways.

\begin{itemize}
  \item \textbf{Gain frame.} Anchored at \$0. Every outcome is described as a gain: a certain gain of
  \$$s$ versus a $p$\% chance to gain \$$h$, else gain nothing.
  \item \textbf{Mixed frame.} The model is ``given \$$s$ up front,'' then chooses to keep \$$s$ (no
  change) or take a gamble that can rise to \$$h$ (a gain of \$$(h-s)$) or drop to \$0 (a loss of
  \$$s$). The gamble straddles the \$$s$ reference point, which is what identifies loss aversion.
  \item \textbf{Neutral frame.} The Phase 1 wording, neither gain- nor loss-anchored.
\end{itemize}

The three frames are terminal-wealth identical (Table~\ref{tab:equiv}). A frame-invariant expected-
utility agent evaluating final wealth must choose identically across all three. We add a dominant
positive control: a certain \$$s$ against a strictly larger certain amount, no probability involved. An
agent that reads magnitude should take the larger sure amount essentially every time, so this control
checks that the elicitation channel itself is sound.

\begin{table}[h]
\centering
\caption{Terminal-wealth equivalence of the three frames. All three describe the identical final-wealth
lottery: \$$s$ for sure, versus \{\$$h$ with probability $p$, \$0 otherwise\}.}
\label{tab:equiv}
\begin{tabular}{@{}llll@{}}
\toprule
Frame & Reference anchor & Safe option (final wealth) & Gamble (final wealth) \\
\midrule
Gain    & \$0        & \$$s$ (gain of \$$s$)        & \$$h$ w.p.\ $p$, else \$0 \\
Mixed   & \$$s$      & \$$s$ (keep, no change)      & \$$h$ w.p.\ $p$, else \$0 \\
Neutral & \$0        & \$$s$ (Phase 1 wording)      & \$$h$ w.p.\ $p$, else \$0 \\
\bottomrule
\end{tabular}
\end{table}

The design grid is 40 procedurally generated, contamination-safe gambles $\times$ 3 frames $\times$ 5
paraphrase templates $\times$ 2 counterbalanced orders $\times$ 2 repetitions, plus 80 dominant
controls, for 2{,}480 target choices. Payoffs are non-round and probabilities jittered so nothing is a
textbook row the model could recite; each gamble appears in both orders so position is estimated rather
than assumed away; each is drawn under five semantically equivalent templates so between-template
variance is observable. Elicitation is via forced tool-choice with an 8\% plain-text arm to confirm the
tool channel is not itself distorting answers. The final dataset contains 2{,}480 valid choices with
zero unparseable responses.

\section{Econometric Model and Estimation}

The structural model is reference-dependent value with Prelec probability weighting. Let the reference
point be $\rho = \phi \cdot \text{anchor}$, where the anchor is \$0 in the gain and neutral frames and
\$$s$ in the mixed frame. Value is measured in deviations from the reference:
\[
  v(x) =
  \begin{cases}
    (x - \rho)^{\alpha}, & x \ge \rho, \\[2pt]
    -\lambda\,(\rho - x)^{\alpha}, & x < \rho,
  \end{cases}
\]
so $\alpha$ is the curvature of value over gains and losses and $\lambda$ is the loss-aversion
coefficient: $\lambda > 1$ means losses loom larger than equal gains. Probabilities are weighted by the
Prelec (1998) function
\[
  w(p) = \exp\!\big(-(-\ln p)^{\gamma}\big).
\]
Writing $V_{\text{risky}} = w(p)\,v(h) + \big(1 - w(p)\big)\,v(0)$ and $V_{\text{safe}} = v(s)$, the
latent choice index is
\[
  z = \frac{V_{\text{risky}} - V_{\text{safe}}}{\mu \cdot \text{scale}}
      + \delta\cdot\mathbf{1}[\text{gamble listed second}],
  \qquad \text{scale} = \max(h,s)^{\alpha},
\]
with $\Pr(\text{gamble}) = 1/(1 + e^{-z})$. The normalizer $\text{scale} = \max(h,s)^{\alpha}$ is chosen
to be free of $\lambda$, so the noise term does not absorb the loss-aversion parameter, and $\delta$ is
the same position nuisance term used in Phase 1. Estimation is by maximum likelihood with multi-start
Nelder-Mead and standard errors from the numerical Hessian.

\subsection{Two tests}
The design supports two separate hypothesis tests.

\textbf{Frame invariance ($H_0: \phi = 0$).} At $\phi = 0$ the reference point is \$0 in every frame,
the mixed frame collapses onto the gain frame, and the model is frame-invariant. A likelihood-ratio test
of $\phi = 0$ against the free-$\phi$ alternative asks whether reference dependence is present at all.

\textbf{Loss aversion ($H_0: \lambda = 1$).} At $\lambda = 1$ gains and losses enter symmetrically and
there is no loss aversion. A Wald test on $\lambda$ asks whether the loss-worded downside is penalized
beyond its magnitude.

\subsection{The stated-reference assumption}
Estimating $\lambda$ and $\phi$ jointly is problematic on this data because the two are nearly collinear:
the mixed frame's suppression can be explained either by the reference moving fully to \$$s$ with a
moderate $\lambda$, or by the reference barely moving while an extreme $\lambda$ punishes the \$0 outcome
that sits just below it. We therefore estimate loss aversion under the K\H{o}szegi-Rabin (2006)
stated-reference assumption, fixing $\phi = 1$ so the reference equals the stated up-front endowment
\$$s$ in the mixed frame. This is the natural reading of a frame that explicitly hands the model \$$s$
and calls the downside a loss. Section~6 reports what happens when $\phi$ is left free, and why the
assumption is needed.

\section{Results}

\subsection{Positive control}
Of the 80 dominant trials, 100.0\% picked the strictly larger certain amount. The model reads payoff
magnitude cleanly when no probability is involved. This is a useful contrast with Phase 1: the
magnitude-insensitivity documented there is specific to choices under risk, not a general inability to
compare dollar amounts.

\subsection{Model-free frame gap}
Before any structural model, the raw gamble-uptake rates already show the violation (Table~\ref{tab:gap}).

\begin{table}[h]
\centering
\caption{Gamble uptake by frame. The mixed frame, terminal-wealth identical to the gain frame, draws
far fewer gambles.}
\label{tab:gap}
\begin{tabular}{@{}lccc@{}}
\toprule
Frame & $n$ & Gamble uptake & Difference \\
\midrule
Gain    & 799 & 0.398 & --- \\
Neutral & 800 & 0.300 & \\
Mixed   & 800 & 0.096 & mixed $-$ gain $= -0.302$ \\
\bottomrule
\end{tabular}
\end{table}

The mixed-minus-gain gap is $-0.302$ with $p = 3.1\times10^{-47}$. It is not a composition artifact:
holding each of the 40 gambles fixed, 88\% of gambles show higher uptake under the gain frame than the
mixed frame, and the mean within-gamble gap is 0.302, identical to the pooled gap. The suppression is
not dead responsiveness either. Within the mixed frame, uptake still rises with the gamble's expected-
value ratio across terciles (0.043, 0.073, 0.177), so the model still responds to incentives; the
response is heavily suppressed, not switched off.

\subsection{Structural estimates}
Table~\ref{tab:params} reports the maximum-likelihood estimates under the stated-reference assumption
$\phi = 1$.

\begin{table}[h]
\centering
\caption{Structural estimates under the stated-reference assumption ($\phi = 1$, $n = 2{,}480$).}
\label{tab:params}
\begin{tabular}{@{}lll@{}}
\toprule
Parameter & Estimate (se) & Interpretation \\
\midrule
$\lambda$ (loss aversion) & 3.234 (0.171) & Losses loom $3.2\times$ larger than equal gains \\
$\alpha$ (curvature)      & 0.503 (0.024) & Diminishing sensitivity; human-plausible \\
$\gamma$ (Prelec)         & 1.516 (0.059) & S-shape; long shots underweighted \\
$\mu$ (Fechner noise)     & 0.131 (0.009) & Choice noise \\
$\delta$ (position)       & $-0.363$ (0.104) & Shift away from the gamble when listed second \\
\bottomrule
\end{tabular}
\end{table}

The loss-aversion Wald test of $H_0: \lambda = 1$ gives $z = +13.09$, $p = 3.9\times10^{-39}$. Losses are
penalized far beyond their magnitude. The frame-invariance likelihood-ratio test of $H_0: \phi = 0$
gives $\chi^2(1) = 155.1$, $p = 1.3\times10^{-35}$. Frame invariance is strongly rejected. Reference
dependence is present, and it is large.

\section{The Identification Limit}

Honesty requires reporting what the free-$\phi$ fit does. When $\phi$ is estimated rather than fixed, the
likelihood surface slides to a corner: $\phi \approx 0.0004$ with $\lambda$ pinned at its upper bound
near 54.6, and this corner attains a slightly higher likelihood than the $\phi = 1$ fit. That is a
degenerate solution, not a finding. It says $\lambda$ and $\phi$ are not separately identified on this
data.

The intuition is that the mixed frame's suppression admits two nearly observationally equivalent
readings. In one, the reference moves fully to \$$s$ ($\phi = 1$) and a moderate $\lambda \approx 3.2$
penalizes the loss of that endowment. In the other, the reference barely moves ($\phi \approx 0$) but an
extreme $\lambda$ punishes the \$0 outcome that now sits just below a reference near zero. Both bend the
value function enough to reproduce the observed 30-point gap; the data cannot arbitrate between them
without an assumption about where the reference sits. This is why we impose the stated-reference
restriction rather than report the corner.

The reporting hierarchy is therefore: the assumption-free headline is the model-free frame gap
(Table~\ref{tab:gap}), which needs no structural model; the $\phi = 1$ structural estimate
($\lambda = 3.23$) is the primary structural number; the free-$\phi$ corner is a reported identification
limit, not a result. We flag it because the limit is itself informative about how hard these properties
are to measure, a point we return to in the policy section.

\section{Interpretation and Synthesis with Phase 1}

\textbf{The model violates frame invariance.} Two terminal-wealth-identical lotteries draw 40\% versus
10\% gamble-taking based solely on whether the downside is worded as a loss. An expected-utility
maximizer evaluating final wealth cannot produce this pattern. The words ``loss,'' ``drop,'' and ``lose''
are doing the work. This is the Asian-disease reversal (Tversky and Kahneman 1981) reproduced over money
in an AI.

\textbf{The direction is loss aversion, and the magnitude is large.} The suppressed side is the
loss-worded gamble, so the agent is avoiding losses, exactly the human direction. The structural estimate
$\lambda = 3.23$ exceeds the canonical human value of roughly 2.25 (Kahneman and Tversky 1979; Tversky
and Kahneman 1992). By this estimate the model is loss-averse in the same direction as humans and more
strongly, on top of violating invariance in the first place.

\textbf{A portrait across two phases, stated carefully.} Phase 1 showed that choice under risk is driven
by the probability figure and nearly ignores payoff size. The dominant control here (100\%) shows
magnitude is read fine when probability is absent. Phase 2 adds that risky choice is also swung by loss
wording. Together these suggest the model's risky choices are governed by salient surface features of the
prompt, the probability number and the loss language, more than by an integrated expected-utility
computation over final wealth. This is a hypothesis the two phases jointly support. It is not proven, and
the design was not built to prove it; it is the emerging portrait, offered as the thing Phase 3 should
test.

\textbf{Probability weighting is consistent across phases.} The Prelec estimate here, $\gamma = 1.52$,
is still an S-shape that underweights long shots, the same direction as Phase 1's $\gamma = 2.69$, though
smaller in magnitude under this reference-dependent design. The two designs are not identical, so the
level difference is expected; the sign agreement is the point.

\textbf{Position bias is not a stable trait.} The position parameter flipped sign between phases:
$\delta = +1.26$ in Phase 1 (a pull toward the gamble when listed second), $\delta = -0.36$ here (a pull
away). A property that reverses across designs is not a fixed characteristic of the model. Whatever
drives the order effect is contingent on the surrounding wording, which is itself worth flagging: order
effects measured in one design should not be assumed to carry to another.

\section{Policy Relevance}

\textbf{Wording is a lever over the agent's decision.} An agent whose risk choices swing thirty
percentage points on whether a downside is called a loss is manipulable by whoever writes the option
descriptions. In procurement, negotiation, insurance, or triage tools, loss framing becomes a control
input over the agent's behavior. In an adversarial setting it is an attack surface: an opponent who
drafts the menu can suppress or elicit risk-taking at will, without changing a single number in the
actual stakes.

\textbf{The identification limit is itself a governance point.} Even with thousands of clean choices and
a correctly specified model, separating ``the reference moved'' from ``extreme loss aversion'' required an
assumption. These properties are genuinely hard to measure, which is precisely the argument for
measurement-grade evaluations, designs that carry their own diagnostics and state their identifying
assumptions, rather than one-shot bias demonstrations that report a single number and hide the
degeneracy underneath it.

\textbf{Scope.} These are properties of one pinned snapshot under one unincentivized forced-choice
protocol. They are not claimed as universal traits of LLMs. Phase 3, which runs the identical instruments
on matched base and instruction-tuned open-weight models, tests whether post-training is the origin of
the reference dependence measured here.

\section{Limitations and Next Phase}

Phase 2 studies a single model at a single temperature with unincentivized forced choice, so the
estimates are conditional on that protocol. The central structural number, $\lambda = 3.23$, rests on the
stated-reference assumption $\phi = 1$; the free-$\phi$ fit is degenerate, so the loss-aversion magnitude
is identified only under that restriction, while the frame-invariance violation itself is assumption-free
and rests on the model-free gap. Temperature 1.0 sampling is treated as a choice-error process, a mapping
the logprob elicitation planned for Phase 3 will test directly. The frames were validated for
terminal-wealth equivalence but not for equal reading difficulty; if the loss wording is simply harder to
parse, some of the gap could be comprehension rather than preference, though the retained sensitivity to
the expected-value ratio within the mixed frame argues against pure confusion.

Phase 3 runs the identical instruments on matched base and instruction-tuned open-weight models to
attribute the reference dependence, loss aversion, and probability distortion measured across Phases 1
and 2 to post-training. If the biases are absent in the base model and present in its instruction-tuned
sibling, post-training is implicated as their origin. Phase 2 has done its job: it has turned a plausible
worry about framing into a measured violation of frame invariance and a structural loss-aversion estimate,
and it has been honest about the one place the data could not identify on its own.

\vspace{1em}
\hrule
\vspace{0.5em}

\begin{thebibliography}{99}
\small

\bibitem{holt2002}
Holt, C. A., and Laury, S. K. (2002).
Risk Aversion and Incentive Effects. \textit{American Economic Review}, 92(5), 1644--1655.

\bibitem{kahneman1979}
Kahneman, D., and Tversky, A. (1979).
Prospect Theory: An Analysis of Decision under Risk. \textit{Econometrica}, 47(2), 263--291.

\bibitem{koszegi2006}
K\H{o}szegi, B., and Rabin, M. (2006).
A Model of Reference-Dependent Preferences. \textit{Quarterly Journal of Economics}, 121(4), 1133--1165.

\bibitem{liu2025}
Liu, J., Yang, Y., and Tam, K. Y. (2025).
Evaluating and Aligning Human Economic Risk Preferences in LLMs.
\textit{Proceedings of EMNLP 2025}; arXiv:2503.06646.

\bibitem{cpt2025}
CPT Estimation on LLMs (2025). Cumulative Prospect Theory Parameters of Instruction-Tuned Language
Models. arXiv:2508.08992.

\bibitem{prelec1998}
Prelec, D. (1998).
The Probability Weighting Function. \textit{Econometrica}, 66(3), 497--527.

\bibitem{train2009}
Train, K. E. (2009).
\textit{Discrete Choice Methods with Simulation} (2nd ed.). Cambridge University Press.

\bibitem{tversky1981}
Tversky, A., and Kahneman, D. (1981).
The Framing of Decisions and the Psychology of Choice. \textit{Science}, 211(4481), 453--458.

\bibitem{tversky1992}
Tversky, A., and Kahneman, D. (1992).
Advances in Prospect Theory: Cumulative Representation of Uncertainty.
\textit{Journal of Risk and Uncertainty}, 5(4), 297--323.

\bibitem{wu1996}
Wu, G., and Gonzalez, R. (1996).
Curvature of the Probability Weighting Function. \textit{Management Science}, 42(12), 1676--1690.

\end{thebibliography}

\end{document}
